\section{Introduction}

The previous chapter established that CI repair is already a recognized research task, but that the most
defensible remaining gap lies in \emph{safe}, \emph{abstention-aware}, and \emph{operationally disciplined}
automation. This chapter turns that positioning into a precise problem statement. Rather than treating PatchOps
as a generic coding agent, the thesis defines it as a constrained decision system for failed GitHub Actions
workflows. The system is designed to observe a failure, compress the relevant evidence, reason over repository
context, and then choose an appropriate operational response under explicit risk constraints.

This precision matters for two reasons. First, it prevents overclaiming. A thesis that vaguely states ``PatchOps
fixes CI failures'' would blur together diagnosis, patch generation, workflow validation, and safety. Second, a
well-defined problem statement makes the evaluation meaningful. The benchmarks later in the thesis measure not
only whether a patch can be produced, but whether the system chooses the right category of action for the right
kind of failure.

\section{Task Definition}

PatchOps addresses the following task: given a failed GitHub Actions workflow execution and, when available, a
local clone of the associated repository, determine the most appropriate response among a bounded set of
operational actions and support that action with traceable evidence. The task is therefore not merely
``generate a patch.'' It is a structured decision problem over noisy evidence.

At a high level, the system input consists of:

\begin{enumerate}
\item the raw CI log or failure trace associated with a failed workflow run,
\item optional repository context, including workflow files, manifests, and likely fix files, and
\item optional execution context used for downstream validation, such as a local reproducible runner.
\end{enumerate}

The system output is a structured result containing compressed evidence, a failure classification, a root-cause
summary, a final decision, and when appropriate a candidate patch and its safety assessment. In the current
implementation this structure is explicit in the engine types: evidence, classification, root cause, patch,
safety report, and decision are all materialized rather than left implicit in free-form text. This is an
important design choice because it makes the system measurable and auditable.

\section{Decision Space}

The central abstraction in PatchOps is the \emph{decision} rather than the patch. The system currently supports
four explicit actions:

\begin{enumerate}
\item \textbf{Patch:} generate a patch that is eligible for automatic progression because the failure is judged
low risk, confidence is sufficiently high, and the patch passes the safety and plausibility checks.
\item \textbf{Preview:} generate a candidate patch or repair suggestion, but keep a human in the loop rather
than treating the change as ready for automatic action.
\item \textbf{Rerun:} recommend re-execution when the failure is likely transient or flaky rather than a stable
code or configuration defect.
\item \textbf{Escalate:} decline automatic repair and forward the case for human review, typically because the
failure is high risk, semantically unclear, or outside the safe editing envelope.
\end{enumerate}

The broader thesis uses the language of \emph{abstention}, and that remains conceptually correct. In PatchOps,
abstention is realized operationally through the non-patch branches of the decision space. A rerun is an
abstention from editing; an escalate decision is an abstention from automated repair; and preview is a
human-in-the-loop abstaining mode for cases where a suggestion may be useful but an autonomous patch would be
too strong. This interpretation keeps the thesis consistent with the literature while matching the implemented
engine state machine exactly.

\section{Failure Taxonomy and Risk Framing}

PatchOps uses a fixed operational taxonomy of failure categories, because the action policy depends on the type
of failure being handled. The current taxonomy contains ten core categories, with an \texttt{unknown} fallback
for unresolved cases:

\begin{center}
\small
\begin{tabularx}{0.92\textwidth}{
  >{\ttfamily\raggedright\arraybackslash}X
  >{\ttfamily\raggedright\arraybackslash}X
}
dependency\_install\_failure & unit\_test\_failure \\
runtime\_version\_mismatch & missing\_system\_dependency \\
workflow\_yaml\_error & permission\_or\_secret\_failure \\
docker\_build\_failure & flaky\_or\_transient\_failure \\
lint\_format\_failure & lockfile\_mismatch \\
\end{tabularx}
\end{center}

These categories are not arbitrary labels invented only for implementation convenience. They are the project's
operational abstraction over three sources: the broader GitHub Actions failure literature, the twelve-category
structure of CI-Repair-Bench, and the practical needs of a decision policy that must distinguish repairable
mechanical failures from risky or ambiguous ones. In later chapters, this taxonomy will be mapped explicitly to
the literature-backed crosswalk used throughout the evaluation.

Each category is also assigned a risk class. In the present system, low-risk categories include formatting,
version mismatches, lockfile issues, dependency installation issues, and workflow YAML fixes. Medium-risk
categories include Docker build failures, unit test failures, and missing system dependencies. High-risk cases
include permission and secret failures. Transient failures are handled separately through rerun logic, while
unknown cases are treated conservatively. This risk structure is one of the mechanisms that turns PatchOps from
a patch generator into a risk-aware decision system.

\section{Confidence, Learned Decision Gate, and Intended Actions}

The decision policy in PatchOps is not purely symbolic; it combines category-level risk with confidence-based
decision logic. In the earlier implementation, the engine used two pinned thresholds: a lower confidence floor
of 0.50 and a higher patch threshold of 0.75. These values made the policy reproducible, but they were still
hand-set constants. The later Phase 15 refinement replaces the low-risk act-versus-hold test with an optional
learned, isotonic-calibrated decision gate over pre-generation signals such as calibrated classifier
confidence, category risk, flaky status, semantic ambiguity, log-named-file count, and evidence length.

Formally, let \(x_i\) denote the pre-generation feature vector for instance \(i\), and let
\(x_i' = (x_i - \mu) / \sigma\) denote the standardized feature vector used by the decision model. The
classifier-facing confidence signal used throughout the policy is
\begin{equation}
\mathrm{conf}_i = \max_c \Pr(c \mid x_i) = \max_c \mathrm{softmax}(W\phi(\mathrm{evidence}_i))_c,
\end{equation}
where \(c\) ranges over failure categories and \(\phi(\cdot)\) denotes the evidence representation. On the
low-risk lane, Phase 15 adds a learned act-versus-hold model:
\begin{equation}
\Pr(\mathrm{act} \mid x_i) = \sigma(w^\top x_i' + b),
\qquad
\sigma(z) = \frac{1}{1 + e^{-z}}.
\end{equation}

Because raw probabilities can be miscalibrated, PatchOps applies a non-decreasing isotonic calibrator
\(g(\cdot)\), fit on labeled act/hold outcomes:
\begin{equation}
g^\star = \arg\min_{g \ \mathrm{monotone}} \sum_i \left(g(s_i) - y_i\right)^2,
\qquad
p_{i,\mathrm{cal}} = g^\star\!\left(\Pr(\mathrm{act} \mid x_i)\right).
\end{equation}
The final low-risk decision rule is then
\begin{equation}
d_i =
\begin{cases}
\textsc{act} & \text{if } p_{i,\mathrm{cal}} \ge \tau, \\
\textsc{hold} & \text{if } p_{i,\mathrm{cal}} < \tau,
\end{cases}
\qquad \tau = 0.5.
\end{equation}
In plain terms, the gate predicts whether a low-risk case should proceed toward an automatic patch or remain in
preview mode.

The intended action is computed before patch generation. If a failure is classified as flaky or transient, the
system chooses rerun rather than edit. If the category is high risk or the diagnosis confidence falls below the
lower floor, the system escalates rather than guessing. Medium-risk categories default to preview. Low-risk
categories become automatic patch candidates only when the confidence policy supports action. In the threshold
baseline this means crossing the higher threshold; in the learned-gate version this means the calibrated
decision model predicts act rather than hold. Otherwise the case remains preview-only.

The older hand-set rule can therefore be written explicitly as
\begin{equation}
\mathrm{conf}_i \ge \mathrm{TAU\_HIGH} = 0.75.
\end{equation}
This mathematical contrast makes the ``remove the magic constant'' point concrete: the learned gate replaces a
single pinned threshold with a calibrated decision surface over multiple pre-generation signals. At the same
time, the broader routing policy remains intentionally deterministic outside this low-risk lane: flaky cases
still rerun, high-risk or unknown cases still escalate, and medium-risk cases still default to preview.

This design reflects the thesis claim that the cost of a wrong action is asymmetric. A false negative, where
the system declines to patch a fixable low-risk issue, is inconvenient. A false positive, where the system
opens or endorses an unsafe or semantically misguided patch, is often more damaging. The policy therefore
intentionally biases the system toward caution.

\section{Semantic Ambiguity and Human Judgement}

One of the most important requirements in PatchOps is that not every technically detectable error should be
treated as machine-fixable. Some failures require genuine human judgement about intended program behavior. The
current implementation encodes this idea through a semantic ambiguity backstop that catches signals such as type
errors, assertion mismatches, and logic-oriented exceptions. Even if such a case appears low risk at the raw
category level, the system downgrades an intended automatic patch to preview.

This requirement is conceptually important because it separates mechanical repair from semantic repair. A format
or version mismatch usually has a narrow and observable fix space. A type mismatch, failed assertion, or logic
exception may involve the intended behavior of the software, which is something the thesis does not claim to
solve fully. By defining semantic ambiguity as a first-class boundary, PatchOps makes its limits explicit rather
than leaving them implicit in later error analysis.

\section{Safety Requirements}

The core safety requirement of PatchOps is straightforward: the system must never treat patch generation as a
license to edit arbitrarily. A candidate patch must pass deterministic safety checks before it can remain in the
\texttt{patch} branch of the decision pipeline.

The current safety envelope has several parts:

\begin{enumerate}
\item \textbf{Forbidden-path protection.} Certain paths, such as environment files, secret-bearing locations,
credentials, and private-key material, are never valid edit targets.
\item \textbf{Category-specific allow-lists.} For categories such as workflow YAML errors, runtime version
mismatches, lockfile issues, and dependency installation failures, the patch must stay within a narrow set of
allowed file patterns.
\item \textbf{Secret-introduction checks.} Added lines are scanned for patterns that resemble API keys,
passwords, tokens, or private keys.
\item \textbf{Minimality constraints.} Automatic patches are bounded by maximum file count and maximum added-line
count, reducing the chance of broad or unfocused edits.
\item \textbf{Reward-hacking detection.} The system checks for high-severity patterns such as test tampering,
assertion removal, skip-based evasion, lint suppression, or CI-disabling changes that could make a run green
without constituting an honest fix.
\end{enumerate}

These requirements embody the thesis claim that safety must be engineered explicitly. They also ensure that
PatchOps can be evaluated on more than repair success alone. Because the safety logic is encoded deterministically,
the later unsafe-edit and abstention metrics are grounded in observable policy rather than vague intuition.

\section{Problem Boundaries and Scope}

The thesis scope is deliberately narrower than the full product fantasy of an autonomous DevOps engineer. The
core comparable evaluation is Python-first because CI-Repair-Bench, the main public benchmark used throughout
the project, is Python-only. Node and Docker examples remain useful in demonstrations and controlled cases, but
they are not mixed into the main comparable repair table.

Similarly, PatchOps does not claim to emulate all of GitHub Actions locally. Its validation harness is a
restricted CI-style re-execution loop meant for supported workflow shapes rather than a perfect clone of every
GitHub-hosted runner feature. This limitation is not a weakness to hide. It is part of the problem definition:
the thesis studies whether a safety-gated agent can make better decisions under a realistic but bounded
validation regime.

The project also does not claim to solve semantic correctness in the strong sense. CI-green is treated as
necessary but not sufficient evidence, and the system is explicitly designed to abstain, preview, or escalate
when the evidence suggests that automated repair would be overconfident. This stance is consistent with both the
APR overfitting literature and the measured behavior of PatchOps in later chapters.

\section{Non-Goals}

Several tempting claims are explicitly outside the scope of this thesis.

\begin{enumerate}
\item \textbf{State-of-the-art repair leadership} is not the primary claim. PatchOps is not framed as a system
whose main contribution is beating the best published Pass@1 number.
\item \textbf{Universal workflow support} is not claimed. The thesis does not assert support for every language,
every CI platform, or every GitHub Actions feature.
\item \textbf{Fully autonomous software engineering} is not claimed. PatchOps is a bounded CI-repair decision
system, not a general autonomous development agent.
\item \textbf{Semantic correctness guarantees} are not claimed. A validated patch is stronger evidence than a
free-form suggestion, but it is still not treated as a perfect proof of software correctness.
\item \textbf{Model-training novelty} is not the center of the thesis. Although the project includes a small
learned classifier as a reproducibility and cost contribution, the main research focus remains orchestration,
safety, abstention, and evaluation.
\end{enumerate}

These non-goals are important because they protect the thesis from being judged against claims it does not need
to make. A well-scoped thesis is stronger than an ambitious but indefensible one.

\section{From Problem Statement to Architecture}

The formal problem statement in this chapter now gives a clean lens for reading the rest of the thesis. PatchOps
accepts failed CI evidence, reasons over bounded repository context, and chooses among patch, preview, rerun,
and escalate actions under explicit risk and safety constraints. The system is therefore defined not by raw
patch generation alone, but by a disciplined pipeline that links diagnosis, action selection, and validation.

This framing also clarifies why the architecture in the next chapter looks the way it does. The multi-stage
pipeline, the safety gate, the semantic ambiguity backstop, the plausibility filter, and the validation loop are
not independent engineering tricks. They are architectural consequences of the problem definition established
here. The next chapter presents that architecture in detail.

\section{Chapter Conclusion}

This chapter has defined the PatchOps task, action space, failure taxonomy, decision policy, safety
requirements, and scope boundaries. These definitions make the later architecture and evaluation measurable:
PatchOps is judged not only by whether it can emit a patch, but by whether it selects an appropriate operational
response under explicit risk constraints. The next chapter translates these requirements into the system
architecture.
